\section{Background}
Biomethane production is going to cover the 4\% of total energy consumption (TEC) of EU by 2050, due to the incentive schemes of the New Green Deal that pushes the energy mix toward low-carbon, renewable and self-produced sources. In the next decade, anaerobic digestion biogas facilities, usually equipped with combined heat\&power engines for the main goal of electricity production, are going to be \emph{upgraded} to biomethane production plants; the produced energy carrier can thus be directly sent to the national grid for third use or stocked in liquid form for transportation use. 

Since biomethane is easier to be stocked than electricity, the focus of plant managers is not going to track a market-driven or demand-driven production, but rather to maximize it. 
As a matter of fact, the incentive scheme (implemented in Italy by "\emph{DM biometano 15 settembre 2022}") has promoted priority of dispatch and fixed tariffs until 2030 (110 \texteuro/MWh, almost double the average methane price of normal market conditions).
Nevertheless, finding the best techno-economic (TE) trade-off point of operation is all but a trivial task, considering also the constraints enforced by the regulating authority to warrant the access to the incentives. 
First of all, co-digestion of different organic substrates (es. energy crops, zootechnical wastes, organic fraction of urban waste) out-competes mono-digestion in terms of performances, but the degree of complexity in finding the best TE operating point increases substantially.

For a successful operation, the plant managers shall be also able to move within the recently-growing agro-industrial by-product (\emph{upstream}) market, where co-substrates are purchased and which rules are all but well-established. 

The anaerobic co-digestion (AcoD) process entails very non-linear and sometimes non-smooth kinetics; in addition, to have a comprehensive and robust "high-fidelity" model able to grasp the behaviour of different complex (es. energy crops and slurries) co-substrates, the state-of-the-art modelling practice requires more than 40 state variables (IWA anaerobic digestion model n.1 (ADM1)).
For this reason, the profit maximization of an already existing plant, may entail a two-stage process: i) a (robust) linear programming problem is solved to find the best influent diet of the reactor; ii) the feasibility of the solution found so far (along with its related worst-case scenario) is verified with the "high-fideliy" non-linear model.

The plant reactor is modelled as a completely stirred tank reactor (CSTR), which foreseen an input flowrate with a given content of organic matter and outputs a gaseous flowrate ($\sim$ 55\% methane and 45\% biogenic carbon dioxide) and a heterogeneous liquid-solid mixture outflow rate called digestate (with its residual unconverted organic matter content). The \emph{biogas} flowrate is thus sent to physico-chemical units to be \emph{upgraded} to biomethane before compression or injection into the grid.
The digestate can be legislatively considered either a by-product or a waste, also depending on the incoming reactor diet. It can thus represent either an additional operational cost for disposal or, when further treated, an additional revenue, since it can be sold to the fertilizers' market, with pricing that depends on its soil nutritional content.  

